% !TEX program = xelatex
\documentclass{英语试卷}
\begin{document}
\examtitle{2025 届「皖南八校」高三第二次大联考}{英语}
\par 英语学科 试题卷\par
\begin{notice}
\par 1. 答卷前，考生务必将自己的姓名、准考证号填写在答题卡上。\par
\par 2. 回答选择题时，选出每小题答案后，用铅笔把答题卡上对应题目的答案标号涂黑。如需改动，用橡皮擦干净后，再选涂其他答案标号。回答非选择题时，将答案写在答题卡上，写在本试卷上无效，\par
\par 3. 考试结束后，将本试卷和答题卡一并交回。\par
\end{notice}
\parttitle{第一部分 听力（共两节，满分 30 分）}
\par 做题时，先将答案标在试卷上，录音内容结束后，你将有两分钟的时间将试卷上的答案转涂到答题卡上。\par
\sectitle{第一节（共5小题：每小题1.5分，满分7.5分）微信公众号IAI English}
\par 听下面 5 段对话。每段对话后有一个小题，从题中所给的 A、B、C 三个选项中选出最佳选项。听完每段对话后，你都有 10 秒钟的时间来回答有关小题和阅读下一小题。每段对话仅读一遍。\par
\question{1. When did the man meet his friend?}
\begin{choices}(3)
  \item At 3:30.
  \item At 4:30.
  \item At 10:30.
\end{choices}
\question{2. How did the woman come to the conference?}
\begin{choices}(2)
  \item By bus. 3. By taxi.
  \item By subway.
\end{choices}
\question{3. Where does the conversation mostly take place?}
\begin{choices}(3)
  \item At the lab.
  \item At the library.
  \item At the classroom.
\end{choices}
\question{4. Why does the son go on a theater trip?}
\begin{choices}(3)
  \item To learn how to put on a play.
  \item To raise money for his club.
  \item To volunteer for the theater.
\end{choices}
\question{5. What are the speakers mainly talking about?}
\begin{choices}(3)
  \item A holiday.
  \item A dish.
  \item A tradition.
\end{choices}
\sectitle{第二节（共15小题：每小题1.5分，满分22.5分）}
\par 听下面 5 段对话或独白。每段对话或独白后有几个小题，从题中所给的 A、B、C 三个选项中选出最佳选项。听每段对话或独白前，你将有时间阅读各个小题，每小题 5 秒钟；听完后，各小题将给出 5 秒钟的作答时间。每段对话或独白读两遍。\par
\par 听下面一段对话，回答第 6 和第 7 两个小题，\par
\question{6. What will the man do on Saturday?}
\begin{choices}(3)
  \item Attend a dancing club.
  \item Organize the audience.
  \item Watch the dance competition.
\end{choices}
\question{7. How does the woman feel about the competition?}
\begin{choices}(3)
  \item Nervous.
  \item Confident.
  \item Calm.
\end{choices}
\par 听下面一段对话，回答第8和第9两个小题。\par
\question{8. What does the woman like best about the song?}
\begin{choices}(3)
  \item Its melody.
  \item Its beats.
  \item Its lines.
\end{choices}
\question{9. What are Eastern Blue's songs usually like?}
\begin{choices}(3)
  \item Shallow.
  \item Meaningful.
  \item Dull.
\end{choices}
\par 听下面一段对话，回答第 10 至第 12 三个小题。\par
\question{10. What does the woman think of the catalogue?}
\begin{choices}(3)
  \item It is useless.
  \item It is detailed.
  \item It is simple.
\end{choices}
\question{11. What is the man satisfied with about the exhibition?}
\begin{choices}(3)
  \item Its use of natural light.
  \item Its supply of energy drinks.
  \item Its gift of free catalogues.
\end{choices}
\question{12. What surprised the man a lot?}
\begin{choices}(3)
  \item The artworks were disappointing.
  \item Few visitors came to the exhibition.
  \item The weather was pleasant.
\end{choices}
\par 听下面一段对话，回答第 13 至第 16 四个小题，\par
\question{13. What is the woman using ChatGPT to do?}
\begin{choices}(3)
  \item Organize outlines.
  \item Generate contents.
  \item Get unique opinions.
\end{choices}
\question{14. What does ChatGPT prefer in the study?}
\begin{choices}(3)
  \item Science fiction.
  \item Romance novels.
  \item Historical books.
\end{choices}
\question{15. How was ChatGPT tested in the study?}
\begin{choices}(3)
  \item By writing sci-fi stories.
  \item By participating in a quiz.
  \item By acting as a student.
\end{choices}
\question{16. Who likes Interstellar most?}
\begin{choices}(3)
  \item Jamie.
  \item Devin.
  \item Dr. Robert.
\end{choices}
\par 听下面一段独白，回答第 17 至第 20 四个小题，\par
\question{17. What has a history of 100 years?}
\begin{choices}(3)
  \item The stone walls.
  \item The iron bridge
  \item The wooden bridge.
\end{choices}
\question{18. What knowledge did the speaker use for the project?}
\begin{choices}(3)
  \item Biology.
  \item Chemistry.
  \item Geography.
\end{choices}
\question{19. Where did the speaker collect water for the project?}
\begin{choices}(3)
  \item A river.
  \item A stream.
  \item A pond.
\end{choices}
\question{20. What is the best presentation about?}
\begin{choices}(3)
  \item The sheep farming.
  \item The bee farming.
  \item The diary farming.
\end{choices}
\parttitle{第二部分 阅读理解（共两节，满分 50 分）}
\sectitle{第一节（共15小题：每小题2.5分，满分37.5分）}
\par 阅读下列短文，从每题所给的 A、B、C 和 D 四个选项中，选出最佳答案，\par
\passagelabel{A}
\par For travelers planning a culture-focused adventure, cities are the best place to start. Some cities function as living galleries in themselves through city-wide street art or award-winning architecture. Others offer artistic ways to explore history and culture, from immersive(沉浸式的) exhibitions to forward-thinking initiatives\par
\par Budapest\par
\par Art takes to the streets in the Hungarian capital. Here, city walls act as a canvas for many homegrown and international artists to make their mark, whether through political messages, nurals(壁画) or admiration for beloved figures and concepts. Though the art form has now spread well across the city, the Rubik's cube and Refugee Girl in the Jewish Quarter remain representative works to get you started.\par
\par Milan\par
\par Milan's fashion heritage dates back to its days as an important trading centre, when rich traders would bring luxurious fabrics to the city and develop and improve their business skills in its streets. It's easy to plan an handcrafted tour taking in contemporary fashion houses, modern art galleries, historic churches and architecture as you go.\par
\par Copenhagen\par
\par Design is integrated into the outline of the Danish capital. Copenhagen's architecture has been awarded many honours over the years, including UNESCO World Capital of Architecture2023. Danish symbol, Arne Jacobsen, was a pioneer of modernist architecture in the city. The spiral(螺旋形的) staircases, Egg and Swan chairs and representative Room 606 at the Radisson Collection Royal Hotel, Copenhagen, were designed by him.\par
\par Riyadh\par
\par Saudi Arabia's largest city has actively positioned itself as a global destination for creative travellers and to this end, it has established the Riyadh Art initiative. The first public art project of its kind in the country, Riyadh Art, aims to transform the city into a gallery without walls through public art festivals and workshops.\par
\par Visit the capital of Saudi Arabian, Riyadh, in winter and you'll see the city shine brightly under the lights of Noor Riyadh, the world's largest light art festival.\par
\question{21. What can visitors do during the tour to Budapest?}
\begin{choices}(2)
  \item Appreciate colorful pictures on the city walls.
  \item Send political messages to beloved figures.
  \item Create representative works across the city.
  \item Spread the art form in their hometown.
\end{choices}
\question{22. What do Copenhagen and Riyadh have in common?}
\begin{choices}(2)
  \item They are capital cities.
  \item They both won awards.
  \item They have galleries without walls
  \item They are famous for ancient architecture.
\end{choices}
\question{23. Who are most likely to show interest in the article?}
\begin{choices}(2)
  \item Traders who sell luxurious fabrics to the city.
  \item Historians who study contemporary history.
  \item Librarians who are interested in literature.
  \item Tourists who prefer a trip focused on culture.
\end{choices}
\passagelabel{B}
\par On a wet June morning, Emily Bick winds through a field of knee-high corn, just north of Madison, Wisconsin. It feels like that quiet, anticipation-filled moment before a concert: Tech people are setting up microphones, laying wires.\par
\par She's here for a show. The star of this particular show is the microphone itself. Research assistants are attaching it to the corn plants, an innovation that Bick named the Insect Eavesdropper(偷听者).\par
\par Bick, a scientist, researches ways to better detect the agricultural insects that lead to serious economic losses worldwide. She says improving the harvest could be done by using chemicals to kill them more strategically — less often, at just the right time.\par
\par In the Corn Belt, that's corn rootworm, also known as the "billion dollar insect". The insect feeds on corn roots, limiting the plant's nutrient supply and making plants easy to falling over. Climate change is bringing new challenges. insects are on the up and up in warmer temperatures and expanding their ranges. "We lose 20\% of our crop to the insect, so we should probably be doing something if we want to conserve more land for non-agricultural use," Bick says.\par
\par Bick's solution to the problem of insect detection is sound: Identify the insect by eavesdropping on the unique shaking that it makes when chewing on a plant. The Insect Eavesdropper uses a contact microphone, an affordable device that musician's often use on instruments. A little metal device sits on a solid object — like corn — and captures the shaking moving through it, registering each one with an electrical signal.\par
\par "When corn rootworm feeds on the roots, the shaking translate from the root system to the top of the plant," Bick says. "That's where we're picking it up. We're using the plant as the musical instrument, as essentially the outside of a guitar."\par
\par The Insect Eavesdropper has opened Bick's ears to a whole world of insect sounds, including her new favorite:\par
\par Aphids, which eat up the sugars flowing through phloem (韧皮部) with a straw-like mouth part. "It kind of sounds like you stuck a straw to the bottom of a milkshake(奶昔)," she says.\par
\question{24. Why does Bick come to the corn field?}
\begin{choices}(2)
  \item To experience the quietness in the wild.
  \item To enjoy a field music performance.
  \item To collect information on agricultural insects.
  \item To anticipate a bright corn harvesting
\end{choices}
\question{25. What can be inferred from Bick's words in Paragraph 4?}
\begin{choices}(2)
  \item The limited supply of corn nutrient.
  \item The urgency to solve the insect attack.
  \item The effect of global warming on agriculture.
  \item The economic loss of the agriculture industry.
\end{choices}
\question{26. How does Bick's Insect Eavesdropper work?}
\begin{choices}(2)
  \item By recording the sound movement from plants.
  \item By broadcasting famous music to corn plants.
  \item By numbering the insects on the plants.
  \item By sending sweet musical signals.
\end{choices}
\question{27. What can we know about Bick's Insect Eavesdropper?}
\begin{choices}(2)
  \item It has high cost.
  \item It has a huge size.
  \item It is reliable.
  \item It is faulty.
\end{choices}
\passagelabel{C}
\par Researchers studied nearly 450 grassy structures built over two years by white-browed sparrow weavers(肯尼亚织雀) living in the Kalahari Desert in South Africa.\par
\par "Experiments show that inexperienced male birds will watch familiar males choosing nest-building materials and then use that information to adjust their own choices. For example, they preferred a pink tape before, rather than an orange one, but they will prefer the orange one if they see others using it," says Maria Cristina Tello-Ramos, a researcher.\par
\par This bird species builds grass structures that serve two purposes: They can be turned into a safe spot to hatch eggs, but they're also shelters where birds can safely rest inside at night. Birds shape grass into a tube, which ends up being shaped basically like an upside- down letter "U".\par
\par To learn more, Maria and her coworkers measured every single structure built by 43 family groups over two seasons. They analyzed everything from the length of the tubes to the tubes' width to the angle between the U-shaped tubes' "arms". They intended to see what might account for the observed differences, carefully considering factors like temperatures, wind speed, and the number of birds in each family. They couldn't find any explanation for why the birds were building such different structures. "And since it's known that these birds are highly social," says Tello-Ramos, "it seems probable that family members are just copying each other, creating a culture of nest-building that gets passed down again and again and again."\par
\par Sometimes a family will take an outsider under their wing and let it join the group. "What I found most fascinating was the fact that when a new sparrow weaver joins the group, it doesn't bring the cultural transmission(传送) of its old group," says Catherine Sheard, a biologist. "That wouldn't have been what I predicted."\par
\question{28. What does Tello stress in her words in the second paragraph?}
\begin{choices}(2)
  \item Birds employ genetics to build nests.
  \item Birds construct nests to attract females.
  \item Birds tend to follow suit in nest-building.
  \item Birds build their nests with colorful materials.
\end{choices}
\question{29. Why do birds build different shapes of nests according to Tello-Ramos?}
\begin{choices}(2)
  \item Because they adapt their sizes to family members.
  \item Because they adjust the nests to match wind speed.
  \item Because they take temperature changes into account.
  \item Because they have various family nesting preferences.
\end{choices}
\question{30. What did Sheard originally expect in her study?}
\begin{choices}(2)
  \item New-comers would build nests according to genetic data.
  \item New birds wouldn't adapt new cultures of nest-building.
  \item New joiners would keep their own cultural traditions.
  \item New members would adopt the new nest-building ways.
\end{choices}
\question{31. What is the best title for the text?}
\begin{choices}(2)
  \item Colorful Nests Are a Fashion for Birds
  \item Tradition Lies in nest-building Materials
  \item Birds Build Their Culture into Their Nests
  \item Nest-Building Reflects Birds' Gender
\end{choices}
\passagelabel{D}
\par Over 100 years ago, Alexander Graham Bell asked the readers to do something brave and fresh — "to found a new science". He pointed out that sciences based on the measurement of sound and light already existed. But there was no science of odor (气味). Bell asked his readers to "measure a smell".\par
\par Today, smart phones provide impressive built-in capabilities based on the sciences of sound and light: voice assistants, facial recognition and photo improvement. The science of odor does not offer anything comparable. But that situation is changing, as advances in machine smell, also called "digitized smell", are finally answering Bell's call to action.\par
\par Whereas human vision mainly relies on cells in the retina(视网膜), smell is experienced through about 400 types of cells in the nose. Machine smell starts with sensors that detect and identify molecules(分子) in the air. These sensors serve the same purpose as the cells in your nose. But to be useful to people, machine smell needs to go a step further. The system needs to know what a certain molecule or a set of molecules smells like to a human. So the challenge is formidable. For that, machine smell needs machine learning.\par
\par Machine learning studies the words humans tend to use to describe what they experience when they encounter specific odor-causing things. However, the web has an unimaginably huge amount of audio, image and video content that can be used to train artificial intelligence systems that recognize sounds and pictures. Most people cannot describe smells in word as effortlessly and recognizably as they can describe sights and sounds. So researchers weren't able to train really powerful machine learning models.\par
\par However, things started to change in 2015 when researchers launched the DREAM Smell Prediction Challenge. The competition released data collected by Andreas Keller and Leslie Vosshall, and invited teams from around the world to submit their machine learning models. By 2019, the largest datasets had grown from less than 500 molecules in the DREAM challenge to about 5,000 molecules. Such progress in unlocking the code of smell is not only exciting but also has highly promising applications, including personalized odor, early detection of disease, and more realistic augmented reality(增强现实技术) experiences. The future of machine smell looks bright. It also promises to smell good.\par
\question{32. Why does the author mention Alexander Graham Bell?}
\begin{choices}(2)
  \item To prove the current progress in machine smell.
  \item To illustrate Bell's research into sound and light.
  \item To stress the application of digitalized smell.
  \item To offer some background information of the science of smell.
\end{choices}
\question{33. What does the underlined word "formidable" in Paragraph 3 mean?}
\begin{choices}(2)
  \item Rewarding.
  \item Demanding.
  \item Stimulating.
  \item Confusing.
\end{choices}
\question{34. What is Paragraph 4 mainly about?}
\begin{choices}(2)
  \item Introduce the shortage of smell datasets.
  \item Warn people of the junk information.
  \item Illustrate the use of smell training model.
  \item Predict the tendency for machine learning.
\end{choices}
\question{35. What is the author's attitude towards the machine smell?}
\begin{choices}(2)
  \item Doubtful.
  \item Optimistic.
  \item Impossible.
  \item Contradictory.
\end{choices}
\sectitle{第二节（共5小题；每小题2.5分，满分12.5分）}
\par 阅读下面短文，从短文后的选项中选出能填入空白处的最佳选项。选项中有两项为多余选项，\par
\par Ever found yourself reaching for your favorite old movie or TV show despite knowing every line and unexpected plot changes? You're not alone. \blanknum{36}\par
\par We rewatch shows for emotion, not plot. Rewatching is an intentional choice based on our prior knowledge. Knowing what's coming allows us to take charge of our feelings. Rewatching and repetition lower our loads of mental understanding too. \blanknum{37} We can freely make ourselves completely involved in the experience, enjoy details we might have missed the first time, and deepen our connection to the characters and storylines.\par
\par \blanknum{38} Returning to something known when life feels chaotic can ground us. Familiar content removes the tension of uncertainty. We can fully focus on the positive moods.\par
\par Rewatching isn't just about the past; it's also about the present and the future. Anticipating a favorite scene plays a significant role. It's reported that the production of dopamine(多巴胺) will be activated by the mental "wanting" of the visualized rewards. \blanknum{39} We get a feel-good boost during the show as we review favorite moments.\par
\par Our media choices reflect our psychological needs. \blanknum{40} So, the next time you queue up old favourites, remember that you're doing more than indulging(沉溺) in nostalgia(怀旧)—you're actively supporting your mental health.\par
\begin{choices}(1)
  \item The familiarity also provides a sense of safety.
  \item Our brains don't need to work hard to process the information.
  \item That's why we simply cannot deal with anything else at the moment.
  \item This mental anticipation motivates behaviors that will deliver the experience.
  \item They can be a response to our desire to recreate meaningful emotional experiences.
  \item Rewatching old favorites is a popular activity for many, especially during the holidays.
  \item A new movie may seem more appealing if we are familiar with the actors, the director or the theme.
\end{choices}
\parttitle{第三部分 语言知识运用（共两节，满分 30 分）}
\sectitle{第一节（共15小题；每小题1分，满分15分）}
\par 阅读下面短文，从短文后各题所给的 A、B、C 和 D 四个选项中，选出可以填入空白处的最佳选项，\par
\par Two years ago, I decided to gather some neighbors to create a community garden. The \blanknum{41} came from two neglected(被忽视的) flower beds at the top of our street, filled with weeds, garbage, and dog wastes. "Why are we \blanknum{42} this?" I thought. I figured anything we did would be an improvement.\par
\par \blanknum{43}, the neighbors were all in. We removed needles, glass, masks, and tons of garbage. We did not receive any funding, so we had to be \blanknum{44}. We started with the smaller bed, getting plants from our gardens. Martha, a chef, \blanknum{45} her roses. "Plant them near the base of the apple tree," she advised. Kishore, with an amazing front garden, \blanknum{46} tulips(郁金香) and carnations(康乃馨). I planted violets. With a few interesting \blanknum{47} well placed, before we knew it, there was a garden.\par
\par The larger bed was still a \blanknum{48}. I would weed it but the soil was \blanknum{49}. When I was thinking about \blanknum{50} it, neighbors across the street had cut branches from a large tree, leaving a giant pile of coverings on the street, so I \blanknum{51} them onto the bed to hold back the \blanknum{52} and to improve the soil.\par
\par "Every time I walk past, I am amazed at how quickly we transform the \blanknum{53}," I said. It's far from \blanknum{54}\par
\cz{41}{depression}{accident}{message}{inspiration}
\cz{42}{wasting}{considering}{wondering}{receiving}
\cz{43}{Strangely}{Frankly}{Luckily}{Apparently}
\cz{44}{serious}{tolerant}{creative}{cautious}
\cz{45}{decorated}{shared}{arranged}{watered}
\cz{46}{contributed}{purchased}{admired}{picked}
\cz{47}{rivers}{mountains}{rocks}{yards}
\cz{48}{mess}{shelter}{trial}{protection}
\cz{49}{high}{heavy}{poor}{sticky}
\cz{50}{taking control of}{giving up on}{holding on to}{making use of}
\cz{51}{burned}{folded}{distributed}{expected}
\cz{52}{flowers}{weeds}{birds}{vehicles}
\cz{53}{space}{power}{strength}{authority}
\cz{54}{welcomed}{changed}{completed}{agreed}
\cz{55}{target}{beauty}{principal}{secret}
\sectitle{第二节（共10小题；每小题1.5分，满分15分）}
\par 阅读下面短文，在空白处填入 1 个适当的单词或括号内单词的正确形式。\par
\par When Chinese developer Game Science revealed its computer game Black Myth(神话故事): Wukong, it immediately gained widespread attention. Inspired by \blanknum{56} great 16 th -century Chinese novel, Journey to the West, the action-packed game featured the monkey Sun Wukong battling sword-wielding (挥动宝剑的) foxes. Smartphone games are \blanknum{57} (extreme) popular in China, \blanknum{58} computer game developers are still few and far between. The enthusiasm for Wukong in Game Science's homeland reached fever pitch(程度) on \blanknum{59} (it) first day. Within 24 hours, the trailer(预告片) achieved 2 million views on YouTube and more than 10 million on Chinese video sharing site Bilibili, much to its creators' shock and joy. One crazy fan even broke into the developer's company, desperate \blanknum{60} more information on the game.\par
\par "We are under great pressure," reflects one of Game Science's co-founders, \blanknum{61} is introduced only as a translator. "This is the first time that we \blanknum{62} (work) on a PC and computer game, and we are also a new studio, so we would really like to thank how \blanknum{63} (tolerate) the players are towards us and their \blanknum{64} (excite) encourages us. When \blanknum{65} (face) with such great encouragement, we try our best to satisfy the expectations of the players and to communicate with players in a very honest way."\par
\parttitle{第四部分 写作（共两节，满分 40 分）}
\sectitle{第一节（满分15分）}
\par 假定你是校学生会主席李华。为弘扬中国优秀传统文化，你校定于下周五举办首届以此为主题的校内英语演讲大赛。请你致函你校外教 Mr. Johnson，邀请他作为本届英语演讲大赛的评委，内容包括：\par
\par (1) 活动目的;\par
\par (2) 活动安排;\par
\par (3) 评判标准;\par
\par (4) 恳请其担任评委， 注意：\par
\par (1) 写作词数应为 80 左右;\par
\par (2) 请按如下格式在答题纸的相应位置作答。\par
\par\vspace{0.4em}
\begin{center}\begin{tabular}{{|l|}}
\hline
Mr. Johnson, \\
 \\
Yours, Li Hua \\
\hline
\end{tabular}\end{center}
\vspace{0.4em}
\sectitle{第二节（满分25分）}
\par 阅读下面材料，根据其内容和所给段落开头语续写两段，使之构成一篇完整的短文，\par
\par In the heart of a city, located between tall buildings and busy streets, was a small but lively animal shelter named "Paws Haven". It had been a place for countless stray(流浪的) and abandoned animals, providing them with food, shelter, and love.\par
\par That day, teenagers Mia and Leo passed the shelter and noticed a poster for a volunteer program.\par
\par Mia was instantly fascinated. "Leo, think about all those poor animals that need our help. Let's sign up to be volunteers." Leo, however, was hesitant. "I don't know, Mia. Volunteering might be a lot of work, and we already have school and other commitments."\par
\par Mia persisted. "Just imagine how rewarding it would be to make a difference in their lives. Come on! Let's do it." Eventually, Leo agreed.\par
\par They went to the animal shelter and were immediately put to work. They cleaned cages, fed the animals, and played with them. As days passed, they grew attached to the furry(覆盖毛皮的) creatures and found a new sense of purpose.\par
\par One sunny Saturday, as they were preparing the afternoon meal for the animals, Mia happened to hear a conversation that filled her with fear. The shelter's director, Ms. Harper, was discussing with a staff member the terrible financial situation Paws Haven was facing. With donations decreasing and operational costs rising, the shelter was on the brink(边缘) of closure. "We can't go on like this much longer," Ms. Harper sighed.\par
\par Mia immediately relayed the news to Leo, and they both felt a sense of urgency. "We can't let this happen," Leo said, determination shining in his eyes. "We have to do something."\par
\par They came up with a plan: organize a community fundraiser event. They would call it "A Night for Paws", a combination of a pet talent show, bake sale, and silent auction(拍卖).\par
\par The following weeks were very busy. Mia managed to secure donations from several local stores, including pet supplies, baked goods, and art pieces for the auction. Meanwhile, Leo created an online page for the event, which quickly gained attention. He also designed flyers and posters, which they distributed around the neighborhood. 注意:\par
\par (1) 续写词数应为 150 左右;\par
\par (2) 请按如下格式在答题卡的相应位置作答，\par
\writingpara{Paragraph 1: The day of the event finally arrived.}
\writinglines{8}
\writingpara{Paragraph 2: Overwhelmed, Ms. Harper; thanked everyone.}
\writinglines{8}
\end{document}
